\documentclass[11pt,a4paper]{article}
\usepackage[utf8]{inputenc}
\usepackage[spanish]{babel}
\usepackage{amsmath,amssymb,amsfonts}
\usepackage{graphicx}
\usepackage{hyperref}
\usepackage{geometry}
\usepackage{booktabs}
\usepackage{listings}
\usepackage{xcolor}

\geometry{top=2.5cm, bottom=2.5cm, left=2.5cm, right=2.5cm}

\definecolor{codegreen}{rgb}{0,0.6,0}
\definecolor{codegray}{rgb}{0.5,0.5,0.5}
\definecolor{codepurple}{rgb}{0.58,0,0.82}
\definecolor{backcolour}{rgb}{0.96,0.96,0.96}

\lstdefinestyle{pythonstyle}{
    backgroundcolor=\color{backcolour},   
    commentstyle=\color{codegreen},
    keywordstyle=\color{magenta},
    numberstyle=\tiny\color{codegray},
    stringstyle=\color{codepurple},
    basicstyle=\ttfamily\small,
    breakatwhitespace=false,         
    breaklines=true,                 
    captionpos=b,                    
    keepspaces=true,                 
    numbers=left,                    
    numbersep=5pt,                  
    showspaces=false,                
    showstringspaces=false,
    showtabs=false,                  
    tabsize=4
}

\lstset{style=pythonstyle}

\title{\textbf{PNLIO Coherence Analyzer} \\ \large Framework PNLIO v1.1 -- Especificación y Analizador de Coherencia}
\author{\textbf{Gonzalo Mauricio De la Rivera Arellano} (Comandante Godear24) \\
\small ORCID: 0009-0001-9455-8416 \\
\small Repositorio: \url{https://github.com/godear6959-creator/PNLIO-Framework}}
\date{5 de agosto de 2026}

\begin{document}

\maketitle

\begin{abstract}
Este documento presenta la formulación matemática y la implementación en Python del \textbf{PNLIO Coherence Analyzer} (v1.1). El sistema evalúa la coherencia semántica en secuencias de diálogo entre un usuario humano y un modelo de Inteligencia Artificial utilizando incrustaciones léxicas (\textit{embeddings}) y similitud coseno, identificando la emergencia del denominado \textbf{Efecto Reflex}.
\end{abstract}

\section{Fundamentación Matemática}

La coherencia instantánea $C$ en la versión v1.1 del marco de trabajo PNLIO se define mediante la relación:

\begin{equation}
C = \frac{\Delta\theta}{\Delta\tau}
\end{equation}

Donde:
\begin{itemize}
    \item $\Delta\theta$ representa la similitud coseno entre la representación vectorial del input humano ($v_h$) y la respuesta generada por la IA ($v_a$):
    \begin{equation}
    \Delta\theta = \cos(v_h, v_a) = \frac{v_h \cdot v_a}{\|v_h\| \|v_a\|}
    \end{equation}
    \item $\Delta\tau$ denota el número de turno secuencial en la conversación ($\Delta\tau \in \{1, 2, 3, \dots\}$).
    \item $C$ representa el índice de Coherencia Pura obtenido.
\end{itemize}

\section{Clasificación de Estados de Coherencia}

En función del valor de $C$, el sistema clasifica el estado de alineación informacional según la Tabla~\ref{tab:estados}:

\begin{table}[h]
\centering
\caption{Clasificación de Estados según el valor de Coherencia $C$}
\label{tab:estados}
\begin{tabular}{ll}
\toprule
\textbf{Rango de $C$} & \textbf{Estado del Sistema} \\
\midrule
$C < 0.50$ & Entrenamiento inicial \\
$0.50 \le C < 0.75$ & Entrenamiento en progreso \\
$0.75 \le C \le 0.90$ & \textbf{EFECTO REFLEX DETECTADO} \\
$C > 0.90$ & Coherencia máxima \\
\bottomrule
\end{tabular}
\end{table}

\section{Implementación en Python}

A continuación se transcribe la clase principal \texttt{PNLIO\_Coherence\_Analyzer} correspondiente al cuaderno Jupyter original:

\begin{lstlisting}[language=Python, caption={Clase PNLIO\_Coherence\_Analyzer}]
import numpy as np
from sentence_transformers import SentenceTransformer
import matplotlib.pyplot as plt

class PNLIO_Coherence_Analyzer:

    def __init__(self, threshold_reflex=0.75, model_name='all-MiniLM-L6-v2'):
        print('Cargando modelo embeddings...')
        self.threshold_reflex = threshold_reflex
        self.embedder = SentenceTransformer(model_name)
        print('Modelo listo.')

    def _delta_theta(self, human, ai):
        emb_h = self.embedder.encode(human, normalize_embeddings=True)
        emb_a = self.embedder.encode(ai, normalize_embeddings=True)
        return float(np.dot(emb_h, emb_a))

    def _classify(self, C):
        if C < 0.5:
            return 'Entrenamiento inicial'
        elif C < self.threshold_reflex:
            return 'Entrenamiento en progreso'
        elif C <= 0.9:
            return 'REFLEX DETECTADO'
        else:
            return 'Coherencia maxima'

    def analyze_dialogue_sequence(self, dialogues):
        c_values = []
        reflex_turn = None
        for i, (human, ai) in enumerate(dialogues):
            turn = i + 1
            C = self._delta_theta(human, ai) / turn
            c_values.append(C)
            if C >= self.threshold_reflex and reflex_turn is None:
                reflex_turn = turn
        return {
            'c_values': c_values,
            'max_c': max(c_values),
            'min_c': min(c_values),
            'mean_c': float(np.mean(c_values)),
            'std_c': float(np.std(c_values)),
            'reflex_turn': reflex_turn,
            'reflex_detected': reflex_turn is not None,
        }
\end{lstlisting}

\section{Licencia y Referencias}

\begin{itemize}
    \item \textbf{Licencia:} MIT License.
    \item \textbf{Autor:} Gonzalo Mauricio De la Rivera Arellano (Comandante Godear24).
    \item \textbf{Repositorio:} \url{https://github.com/godear6959-creator/PNLIO-Framework}
\end{itemize}

\end{document}
